\SourceLeaf{66}
$\beta$) Der auf den Anfangskoeffizienten folgende Koeffizient, den wir $C^{(1)}$ nennen wollen, lässt sich aus der zuvor
angegebenen Rekursionsformel bestimmen. Es gilt
\[
C^{(1)} .  n = -2\int nn\,\mathrm{d}C + nn,
\]
Einsetzen und Ausrechnen ergibt
\[
C^{(1)}n = 2\bigl(1-(1-nn)^{\frac{1}{2}}\bigr),
\]
und somit
\[
C^{(1)} = 2\,\frac{1-\BareRoot{(1-nn)}}{n}.
\]
\par\noindent\hspace*{\parindent}$\gamma$) Die übrigen Koeffizienten lassen sich aus der zuvor hergeleiteten Rekursionsformel durch eine einzige allgemeine
Formel bestimmen. Entwickelt man $\log(1+n\cos\varphi)$ in eine Reihe und bezeichnet den indizierten
Koeffizienten durch $h^{\mathrm{tus}} = C^{(h)}$, so erhält man für ihn
\[
= (-1)^{h+1}\,\frac{2}{h}\left(\frac{1-\sqrt{1-nn}}{n}\right)^{h}.
\]
\par\noindent\hspace*{\parindent}Nehmen wir an, das Gesetz gelte für die ersten $h$ Koeffizienten. Dann ergibt sich aus der allgemeinen Rekursionsformel mit $\nu=0$ nach (7, $\gamma$)
\[
C^{(h+1)} = \frac{-n(h-1)C^{(h-1)}-2h.\,C^{(h)}}{n(h+1)}, 
\]
Setzt man die entsprechenden Werte ein, so ergibt sich
\[
C^{(h+1)} = -\frac{\begin{gathered}n.(h-1)(-1)^{h}\dfrac{2}{h-1}\left(\dfrac{1-\sqrt{1-nn}}{n}\right)^{h-1}\\{}+2h(-1)^{h+1}\dfrac{2}{h}\left(\dfrac{1-\BareRoot{1-nn}}{n}\right)^{h}\end{gathered}}{n.(h+1)},
\]
oder nach Ausklammern des gemeinsamen Faktors
\[
C^{(h+1)} = -\frac{(-1)^{h}2\left(\dfrac{1-\sqrt{1-nn}}{n}\right)^{h-1}}{n(h+1)}\left(n-\frac{2(1-\sqrt{1-nn})}{n}\right).
\]
Der in Klammern eingeschlossene Faktor ist
\[
= \frac{nn-2+2\sqrt{1-nn}}{n} = -\frac{(1-\BareRoot{(1-nn)})^{2}}{n}.
\]
Führt man die Multiplikation mit diesem Faktor aus, so ergibt sich
\[
C^{(h+1)} = (-1)^{h+2}\,\frac{2}{h+1}\left(\frac{1-\BareRoot{(1-nn)}}{n}\right)^{h+1},
\]
woraus die Allgemeingültigkeit des Gesetzes hervorgeht.

